\documentclass[11pt,a4paper,openany]{book}

\makeatletter\def\input@path{{../_shared/}}\makeatother
\usepackage{course}

\begin{document}
\raggedbottom

\coursetitle{Computer Security and Privacy}{BA5}{2026--2027}

\tableofcontents
\clearpage

These notes cover the core vocabulary and reasoning framework of computer security:
adversaries, threat models, security policies, mechanisms, and how to argue that a
system is secure.

% =========================================================
\chapter{Lecture 1 : Basics}
% =========================================================

\begin{rembox}[title=Why a course on security?]
Ordinary engineering asks for \textbf{correctness}, \textbf{safety} and
\textbf{robustness} against accidental errors. Security additionally requires that
properties hold in the presence of a \textbf{deliberate, strategic adversary}.
\end{rembox}

\section{Security properties}

\textbf{Properties} of a computer system must hold in the presence of a \textbf{resourced strategic adversary}.

\begin{defnbox}[title=CIA triad]
\begin{itemize}
  \item \textbf{Confidentiality} — prevention of unauthorized \emph{disclosure} of information.
  \item \textbf{Integrity} — prevention of unauthorized \emph{modification} of information.
  \item \textbf{Availability} — prevention of unauthorized \emph{denial of service}.
\end{itemize}
\end{defnbox}

\begin{exbox}[title=CIA on a bank account]
The adversary should not be able to \textit{read} my statement (confidentiality),
\textit{change} my balance (integrity), or \textit{block} my access to the account
(availability).
\end{exbox}

\begin{rembox}[title=Beyond the CIA triad]
Modern systems also care about \textbf{authenticity}, \textbf{anonymity},
\textbf{isolation}, and more.
\end{rembox}

\section{Security policy}

\begin{defnbox}[title=Security policy]
A high-level description of the \textbf{security properties} that must hold, in
relation to \textbf{assets} and \textbf{principals}.
\begin{itemize}
  \item \textbf{Assets (objects)} — anything of value that needs protection (data, files, memory, \ldots).
  \item \textbf{Principals (subjects)} — people, programs, services, \ldots (never the adversary).
\end{itemize}
\end{defnbox}

\begin{propbox}[title=Properties in terms of principals and assets]
\begin{itemize}
  \item \textbf{Confidentiality} — authorized \textbf{users} may read a \textbf{file}.
  \item \textbf{Integrity} — authorized \textbf{programs} may write a \textbf{file}.
  \item \textbf{Availability} — authorized \textbf{services} can access a \textbf{file}.
\end{itemize}
\end{propbox}

\section{The adversary}

\begin{defnbox}[title=Threat model]
Describes the \textbf{resources} available to the adversary and its \textbf{capabilities}
(observe, influence, corrupt, \ldots).
\end{defnbox}

\begin{defnbox}[title=Adversary]
A malicious entity aiming to breach the security policy. The adversary is
\textbf{strategic}: it always picks the \textbf{optimal} way to use its resources to
violate a security property.
\end{defnbox}

\begin{rembox}[title=Key vocabulary]
\begin{itemize}
  \item \textbf{Vulnerability} — a specific weakness an adversary could exploit (\textit{e.g.\ a password shown in plaintext}).
  \item \textbf{Threat} — the feared event, i.e.\ the adversary's goal (\textit{e.g.\ a hacker wants to steal money from the bank}).
  \item \textbf{Harm} — what actually happens if the threat materializes (\textit{e.g.\ the adversary steals the money}).
\end{itemize}
If the properties hold within the threat model, no vulnerability can be exploited to
realize a threat, so no harm occurs.
\end{rembox}

\section{Security mechanisms}

\begin{defnbox}[title=Security mechanism]
A technical mechanism (software, hardware, cryptography, procedures, \ldots) used to
ensure that the security policy is not violated by an adversary
\underline{within the threat model}.
\end{defnbox}

\begin{exbox}[title=Policy $\to$ mechanism]
\textbf{Policy}: the transaction log must not be tampered with by a single employee. \\
\textbf{Mechanism}: replicate the log across multiple machines so that no employee has
access to all copies.
\end{exbox}

\begin{propbox}[title=Composing mechanisms]
Security does \textbf{not} increase linearly with the number of mechanisms.
\begin{itemize}
  \item \textbf{Defence in depth} (AND) — secure as long as \emph{one} mechanism holds.
  \item \textbf{Weakest link} (OR) — insecure as soon as \emph{any one} mechanism fails.
\end{itemize}
\end{propbox}

\begin{warnbox}[title=Humans are part of the system]
Social engineering, phishing, and weak/reused/written-down passwords are frequently the
weakest link — but that is no excuse to neglect the rest of the system.
\end{warnbox}

\section{Showing a system is secure}

\begin{rembox}[title=Attacker vs.\ defender]
\textbf{Attacker}: one way to violate one property is enough (within the threat model). \\
\textbf{Defender}: must show \textbf{no} adversary strategy violates the policy —
security is always relative to \textbf{a specific threat model}.
\end{rembox}

\begin{defnbox}[title=Secure]
A system is \textbf{secure} if an adversary constrained by a specific threat model
cannot violate the security policy.
\end{defnbox}

\begin{defnbox}[title=Security argument]
A rigorous argument (verbal or mathematical) that the security mechanisms in place are
effective in maintaining the security policy — \textbf{subject to the assumptions of
the threat model}.
\end{defnbox}

\begin{warnbox}[title=Composition is hard]
A useful threat model must constrain the adversary, otherwise no security argument is
possible. Security arguments for the \textbf{composition} of mechanisms are notoriously
hard to get right.
\end{warnbox}

\section{Security engineering process}

\begin{algbox}[title=Three-stage process]
\begin{enumerate}
  \item \textbf{High-level specification} — architecture, security policy (principals, assets, properties), threat model.
  \item \textbf{Security design} — select/design mechanisms; state the security argument (which mechanism maintains which property).
  \item \textbf{Secure implementation} — implement the mechanisms, ensure conformance to the design, perform security testing.
\end{enumerate}
\end{algbox}

\begin{keybox}[title=Summary]
\begin{itemize}
  \item Security problems always involve an \textbf{adversary}, who is \textbf{strategic}.
  \item The adversary's capabilities define a \textbf{threat model}.
  \item \textbf{Security mechanisms} aim at fulfilling a \textbf{security policy within a threat model}.
  \item Showing security means providing a \textbf{security argument}.
\end{itemize}
\end{keybox}

\end{document}
